\section{Related Work}

\subsection{Certificate-Based and Pseudonym Schemes}

The Security Credential Management System (SCMS), standardized under IEEE 1609.2,
provides the baseline security architecture for V2X communications. Vehicles are
issued sets of pseudonym certificates (typically 20--100 per week) and rotate them
periodically to reduce linkability. While SCMS provides a well-understood security
model, it faces two practical challenges. First, the CRL distribution problem: as
the number of revoked pseudonyms grows, CRL sizes increase substantially, and
distributing updated CRLs across a vehicular network with intermittent connectivity
is difficult~\cite{scopelliti2024}. Second, messages signed under the same pseudonym
within a rotation period are linkable, enabling short-term trajectory tracking.

\subsection{Batch Verification Signature Schemes}

Wu et al.~\cite{zhang2024} proposed a certificateless signature scheme with
batch verification for V2V, eliminating certificate management overhead and improving
verification scalability. However, the scheme does not achieve full unlinkability:
signatures still require pseudonymous identifiers linkable within a validity period.

Jing et al.~\cite{zheng2018eaba} proposed the Efficient Anonymous Batch
Authentication (EABA) scheme for VANETs, introducing a message classification
algorithm to prioritize high-importance safety messages and a cooperative
authentication mechanism between RSUs and vehicles. EABA reduces authentication
delay through priority-based parallel processing but relies on pseudonym-based
signatures and does not extend naturally to ZKP-based unlinkable credentials. We
incorporate EABA's priority classification principle into our Receiver Module
(Section~\ref{sec:concurrent}).

\subsection{Zero-Knowledge Proof-Based Vehicular Authentication}

ZKP-based approaches aim to eliminate the linkability problem entirely by allowing
vehicles to prove authorization without revealing any identifier. However, existing
works expose two recurring gaps specific to high-frequency V2V.

Gabay et al.~\cite{gabay2020} proposed a privacy-preserving authentication scheme
for connected \emph{electric vehicles} using blockchain and ZKP group membership
proofs. Their scheme targets session-level authentication for EV charging contexts,
not 10\,Hz broadcast safety message authentication. The blockchain
consensus-based credential verification introduces latency incompatible with
sub-100\,ms safety cycles.

Jiang and Guo~\cite{jiang2024} designed an anonymous authentication scheme for the
Internet of Vehicles using a main-secondary blockchain with ZKP. Pathak et
al.~\cite{pathak2024} presented a blockchain-enhanced ZKP scheme for IoT mutual
authentication. Both achieve strong anonymity but treat each authentication as a
full-cost independent ZKP operation without session-level proof reuse or
precomputation optimization.

Zheng et al.~\cite{zheng2024hybrid} proposed a hybrid message authentication
scheme for IoV combining ZKP for initial credential establishment with ECDSA for
per-message signing. While this reduces per-message overhead, the ECDSA component
reintroduces pseudonymous linking and the scheme does not support batch verification.

Smahi et al.~\cite{smahi2023bvicvs} integrated zkSNARKs into a federated learning
framework for V2X environments (BV-ICVs), demonstrating practical zkSNARK deployment
in vehicular settings. However, their work targets data integrity for model updates
rather than high-frequency safety message authentication.

Liu et al.~\cite{liu2025} proposed a Merkle hash tree-based multi-access
authentication scheme for vehicular networks that combines bilinear pairings with
non-interactive ZKPs for initial credential issuance, and uses the Merkle tree as a
lightweight access-counting structure for subsequent V2I authentications. While
this work demonstrates the effectiveness of Merkle trees for reducing per-access
overhead, it focuses on vehicle-to-\emph{infrastructure} authentication with server
interaction on each cycle and does not address the ultra-low-latency requirements
of direct V2V safety message broadcasting at 10\,Hz intervals without any server
interaction.

\subsection{Batch and Group Authentication in VANETs}

Bagga et al.~\cite{bagga2020bbas} proposed BBAS-IoV combining V2V single
authentication with blockchain-based batch cluster authentication. The BATKA
scheme~\cite{batka2025} proposes ECC-based batch authentication with a tree-based
group key structure for 5G-VANETs, reducing RSU verification overhead by grouping
vehicles into clusters. However, BATKA assumes cluster stability that may not hold
in high-speed highway scenarios: at a 100\,km/h relative lane speed difference, a
vehicle exits a 300\,m cluster radius in approximately 10.8\,s. Furthermore,
neither scheme achieves unlinkability within the batch window.

Our scheme extends the batch verification principle to the Groth16 zk-SNARK setting,
exploiting the algebraic structure of pairing-based proofs to achieve clean $k+3$
pairing batch verification while maintaining full unlinkability---a property not
achieved by any signature-based batch scheme.

\subsection{Channel Impairments and Authentication Resilience}

A limitation of most V2V authentication analyses is evaluation under ideal channel
assumptions. Real-world IEEE 802.11p DSRC channels exhibit packet loss rates (PLR)
of 10--40\% at 150--300\,m range under non-line-of-sight (NLoS) conditions---the
dominant scenario in dense highway traffic~\cite{ramadan2025}. Under vehicle
densities $>$100 vehicles/km, the EDCA mechanism becomes saturated, further
increasing PLR. Ramadan et al.~\cite{ramadan2025} highlight that a cross-layer
perspective integrating authentication overhead with MAC-layer congestion control
is necessary for realistic performance analysis. We address these impairments
through our smaller per-message payload ($\approx$2$\times$ smaller than SCMS)
and batch verification's graceful packet-loss degradation (Section~\ref{sec:noise}).

\subsection{Identified Gaps}

Table~\ref{tab:comparison} summarises the comparison. The key gaps we address are:
(i)~\emph{lack of session-based proof reuse} in existing ZKP schemes;
(ii)~\emph{no offline/online decomposition} for high-frequency ZKP generation;
(iii)~\emph{no Groth16 batch verification} for vehicular ZKP authentication;
(iv)~\emph{insufficient treatment of AST issuance security, request policy, and
epoch management}; and (v)~\emph{absence of real-world channel resilience analysis}.

\begin{table}[t]
\centering
\caption{Comparison of V2V Authentication Approaches}
\label{tab:comparison}
\resizebox{\columnwidth}{!}{%
\begin{tabular}{@{}lcccccc@{}}
\toprule
\textbf{Scheme} & \textbf{Unlink.} & \textbf{No CRL} & \textbf{Precomp.} & \textbf{Batch} & \textbf{V2V} & \textbf{Epoch} \\
\midrule
SCMS/ECDSA & \texttimes & \texttimes & \texttimes & \checkmark & \checkmark & \texttimes \\
Wu et al.~\cite{zhang2024} & \texttimes & \checkmark & \texttimes & \checkmark & \checkmark & \texttimes \\
Gabay et al.~\cite{gabay2020} & \checkmark & \checkmark & \texttimes & \texttimes & \texttimes & \texttimes \\
Jiang \& Guo~\cite{jiang2024} & \checkmark & \checkmark & \texttimes & \texttimes & \texttimes & \texttimes \\
Pathak et al.~\cite{pathak2024} & \checkmark & \checkmark & \texttimes & \texttimes & \texttimes & \texttimes \\
Zheng et al.~\cite{zheng2024hybrid} & \texttimes & \checkmark & \texttimes & \texttimes & \checkmark & \texttimes \\
EABA~\cite{zheng2018eaba} & \texttimes & \texttimes & \texttimes & \checkmark & \checkmark & \texttimes \\
\textbf{ULP-V2V-Auth (Ours)} & \checkmark & \checkmark & \checkmark & \checkmark & \checkmark & \checkmark \\
\bottomrule
\end{tabular}%
}
\end{table}
